\section{重心高さ変動と ZMP の関係}

従来の多くの歩行制御手法は，線形倒立振子モデル（LIPM）に基づき，重心（Center of Mass: CoM）の高さが一定であることを仮定している．この仮定の下では，CoM の水平方向運動と ZMP の関係が単純化され，歩行生成および安定化制御が容易となる．

しかし，軟弱地面において床沈み込みが生じる場合，支持脚側の床変形によりロボット全体の CoM 高さが低下する．この CoM 高さの変動は，ZMP 応答に直接的な影響を与える．同一の水平方向加速度が与えられた場合であっても，CoM 高さが変化することで ZMP の位置が変動し，支持多角形内に ZMP を維持することが困難となる．

特に片脚支持期では支持多角形が小さいため，わずかな CoM 高さ変動であっても ZMP の安定余裕が大きく低下する．このことは，軟弱地面において CoM 高さ変動を無視できないことを示唆している．